\documentclass[10pt,a4paper]{amsart}
\usepackage[margin=19mm]{geometry}
\usepackage{amsmath,amssymb,mathtools}
\usepackage[scr=rsfs,cal=euler]{mathalfa}
\usepackage{xfrac}
\usepackage[unicode,hidelinks,bookmarks=false]{hyperref}
\newcommand{\ie}{i.e.\ }
\newcommand{\cf}{cf.\ }
\newcommand{\resp}{respectively\ }
\newcommand{\Subsection}{\subsection}
\providecommand{\nobreakdash}{\nobreak\hbox{-}}
\setlength{\emergencystretch}{2em}
\multlinegap0pt
\usepackage{amssymb}
\newtheorem{theorem}{Theorem}
\newcommand{\FF}{\mathbb{F}}
\newcommand{\ZZ}{\mathbb{Z}}
\title{Ordinary $3$-Isogeny Graphs and Improvement of Supersingularity Testing for Twisted Hessian Curves over Prime Fields}
\author{Yuji Hashimoto, Koji Nuida}
\date{Source: arXiv:2609.11037v1 (CC BY 4.0)}
\begin{document}
\maketitle

\noindent\emph{Source and licence.} Ordinary $3$-Isogeny Graphs and Improvement of Supersingularity Testing for Twisted Hessian Curves over Prime Fields. Version arXiv:2609.11037v1 (CC BY 4.0), distributed by arXiv under the Creative Commons Attribution 4.0 International licence, http://creativecommons.org/licenses/by/4.0/, as recorded in the arXiv licence metadata for this exact version and shown on the abstract page https://arxiv.org/abs/2609.11037v1. Full licence notice: \texttt{licenses/arxiv-2609.11037-CC-BY-4.0.txt}.\par\smallskip

\noindent\textit{Editorial scope.} Counted unit: the article's theorem thm:whether\_codomain\_is\_over\_Fp2 (source0.tex lines 336--341). The article's complete original proof of it is delivered with it (source0.tex lines 342--606, one \texttt{proof} environment of 265 lines). No mathematics is written, repaired or re-derived: the statement and the proof are the article's own lines, in the article's own notation, and no retained line is substituted or re-flowed.  No further source-local context is needed: every cross-reference of every retained line resolves inside the two delivered spans.  Nothing was omitted from any retained span, so the package carries no omission record.  No retained line contains a citation, so the wrapper carries no bibliography and no bibliography entry.  No retained line contains a figure environment, an image-inclusion command or drawing code, so no image is delivered and the package declares no asset dependency.  Editorial formatting only, in addition: an AMS \texttt{amsart} preamble that restates the article's own declarations and macros used by the retained lines, namely the environments \texttt{theorem} on separate counters, together with the standard packages those lines need.  The retained environments print in this document as Theorem 1 in order of appearance, whereas the article prints the same items with the numbering of its own full document.  The complete native source of the article as distributed in the arXiv e-print for this exact version is bundled unchanged as \texttt{sources/209981/source0.tex}.
\begin{theorem}
\label{thm:whether_codomain_is_over_Fp2}
Let $H(a,d)$ be a twisted Hessian curve with $a,d \in \FF_{p^2}$ and $d \neq 0$.
Let $\phi \colon H(a,d) \to H(A,D)$ be a $3$-isogeny of type $1$ associated with $c \in \overline{\FF_p}$, where $c^3 = a$, $A = d^2 c + 3 d c^2 + 9a$, and $D = d + 6c$.
If $j(H(A,D)) \in \FF_{p^2}$, then $c \in \FF_{p^2}$.
\end{theorem}

\begin{proof}
Assume, for the sake of contradiction, that $c \not\in \FF_{p^2}$.
Let $\omega \in \FF_{p^2}$ be an element with $\omega^2 + \omega + 1 = 0$, which is a primitive third root of unity.
Then $a$ has three cubic roots $c$, $\omega c$, and $\omega^2 c$, none of which lies in $\FF_{p^2}$ (as $c \not\in \FF_{p^2}$ and $\omega \in \FF_{p^2}$).
Therefore $X^3 - a$ is the minimal polynomial of $c$ over $\FF_{p^2}$.
Now there is an automorphism $\sigma$ of the extension field $\FF_{p^2}(c)/\FF_{p^2}$ satisfying $\sigma(c) = \omega c$.

Let $d_0 := d/3$.
Then we have
\[
\begin{split}
j(H(A,D))
= \frac{ D^3 }{ A } \frac{ (D^3 + 216A)^3 }{ (D^3 - 27A)^3 }
&= \frac{ 3^3 (d_0 + 2c)^3 }{ 9 (d_0{}^2 c + d_0 c^2 + c^3) } \frac{ ( 3^3 (d_0 + 2c)^3 + 216 A )^3 }{ ( 3^3 (d_0 + 2c)^3 - 27 A )^3 } \\
&= \frac{ 3 (d_0 + 2c)^3 }{ c (d_0{}^2 + d_0 c + c^2) } \frac{ ( (d_0 + 2c)^3 + 8 A )^3 }{ ( (d_0 + 2c)^3 - A )^3 } \enspace.
\end{split}
\]
As $d_0{}^2 + d_0 c + c^2 = (d_0{}^3 - c^3) / (d_0 - c) = (d_0{}^3 - a) / (d_0 - c)$ (note that $d_0 - c \neq 0$, as $d_0 \in \FF_{p^2}$ and $c \not\in \FF_{p^2}$) and
\[
\begin{split}
(d_0 + 2c)^3 - A
&= (d_0{}^3 + 6 d_0{}^2 c + 12 d_0 c^2 + 8 c^3) - 9 (d_0{}^2 c + d_0 c^2 + c^3) \\
&= d_0{}^3 - 3 d_0{}^2 c + 3 d_0 c^2 - c^3
= (d_0 - c)^3 \enspace,
\end{split}
\]
we have
\[
\begin{split}
j(H(A,D))
= \frac{ 3 (d_0 + 2c)^3 }{ c (d_0{}^3 - a) / (d_0 - c) } \frac{ ( (d_0 + 2c)^3 + 8 A )^3 }{ (d_0 - c)^9 }
= \frac{ 3 (d_0 + 2c)^3 }{ c (d_0{}^3 - a) } \frac{ ( (d_0 + 2c)^3 + 8 A )^3 }{ (d_0 - c)^8 }
\in \FF_{p^2} \enspace,
\end{split}
\]
therefore, as $d_0{}^3 - a \in \FF_{p^2}$, we have
\[
\frac{ (d_0 + 2c)^3 }{ c } \frac{ ( (d_0 + 2c)^3 + 8 A )^3 }{ (d_0 - c)^8 }
\in \FF_{p^2} \enspace.
\]
This element is fixed by the automorphism $\sigma$, therefore we have
\[
\frac{ (d_0 + 2c)^3 }{ c } \frac{ ( (d_0 + 2c)^3 + 8 A )^3 }{ (d_0 - c)^8 }
= \frac{ (d_0 + 2 \omega c)^3 }{ \omega c } \frac{ ( (d_0 + 2 \omega c)^3 + 8 \sigma(A) )^3 }{ (d_0 - \omega c)^8 } \enspace.
\]
This implies that
\[
F
:= (d_0 + 2c)^3 ( (d_0 + 2c)^3 + 8 A )^3 \omega (d_0 - \omega c)^8
- (d_0 + 2 \omega c)^3 ( (d_0 + 2 \omega c)^3 + 8 \sigma(A) )^3 (d_0 - c)^8
= 0 \enspace.
\]

Now by expanding $F$ and using the relations $c^3 = a$ and $\omega^2 + \omega + 1 = 0$, it follows from a direct calculation that
\begin{equation}
\label{eq:thm:whether_codomain_is_over_Fp2:1}
\begin{split}
F
&= (F_2 \cdot 2^4 \cdot 3 \cdot 5 \cdot 7 \cdot 13 \cdot (2 \omega + 1) \cdot d_0{}^3 \cdot a^5) \cdot c^2 \\
&\quad + (F_1 \cdot 2 \cdot (\omega + 2) \cdot d_0 \cdot a^6) \cdot c \\
&\quad + (F_0 \cdot (\omega - 1) \cdot d_0{}^2 \cdot a^6) \\
&= 0
\end{split}
\end{equation}
where, by putting $z := d_0{}^3 / a$,
\[
F_2 := z^5 + 2614 z^4 - 32336 z^3 - 433360 z^2 - 57856 z + 41984 \enspace,
\]
\[
F_1 := 4 z^6 + 2790621 z^5 + 427259532 z^4 - 546940856 z^3 - 4822215552 z^2 - 316987392 z + 25907200 \enspace,
\]
\[
F_0 := z^6 + 621168 z^5 - 126479844 z^4 - 4334764688 z^3 - 184901856 z^2 + 4578771456 z + 66734080 \enspace.
\]
Indeed, by putting
\[
\widehat{F}_2
:= F_2 \cdot a^5
= d_0{}^{15} + 2614 d_0{}^{12} a - 32336 d_0{}^9 a^2 - 433360 d_0{}^6 a^3 - 57856 d_0{}^3 a^4 + 41984 a^5 \enspace,
\]
\[
\begin{split}
\widehat{F}_1
:= F_1 \cdot a^6
&= 4 d_0{}^{18} + 2790621 d_0{}^{15} a + 427259532 d_0{}^{12} a^2 - 546940856 d_0{}^9 a^3 \\
&\quad - 4822215552 d_0{}^6 a^4 - 316987392 d_0{}^3 a^5 + 25907200 a^6 \enspace,
\end{split}
\]
\[
\begin{split}
\widehat{F}_0
:= F_0 \cdot a^6
&= d_0{}^{18} + 621168 d_0{}^{15} a - 126479844 d_0{}^{12} a^2 - 4334764688 d_0{}^9 a^3 \\
&\quad - 184901856 d_0{}^6 a^4 + 4578771456 d_0{}^3 a^5 + 66734080 a^6 \enspace,
\end{split}
\]
we have
\begin{equation}
\label{eq:thm:whether_codomain_is_over_Fp2:2}
\begin{split}
F
= (\widehat{F}_2 \cdot 2^4 \cdot 3 \cdot 5 \cdot 7 \cdot 13 \cdot (2 \omega + 1) \cdot d_0{}^3) \cdot c^2 
+ (\widehat{F}_1 \cdot 2 \cdot (\omega + 2) \cdot d_0) \cdot c 
+ (\widehat{F}_0 \cdot (\omega - 1) \cdot d_0{}^2)
\end{split}
\end{equation}
in the quotient ring $\ZZ[a,d_0,\omega,c] / (c^3 - a, \omega^2 + \omega + 1)$ when $a$, $d_0$, $\omega$, and $c$ are temporarily regarded as indeterminates, therefore Eq.\eqref{eq:thm:whether_codomain_is_over_Fp2:2} still holds when taking reduction modulo $p$ and substituting the actual values of $a$, $d_0$, $\omega$, and $c$.
Hence Eq.\eqref{eq:thm:whether_codomain_is_over_Fp2:1} holds in $\overline{\FF_p}$.

Regarding Eq.\eqref{eq:thm:whether_codomain_is_over_Fp2:1}, as $F_0,F_1,F_2,d_0,a,\omega \in \FF_{p^2}$ and $X^3 - a$ is the minimal polynomial of $c$ over $\FF_{p^2}$, it follows that
\[
\begin{cases}
F_2 \cdot 2^4 \cdot 3 \cdot 5 \cdot 7 \cdot 13 \cdot (2 \omega + 1) \cdot d_0{}^3 \cdot a^5 = 0 \enspace,\\
F_1 \cdot 2 \cdot (\omega + 2) \cdot d_0 \cdot a^6 = 0 \enspace,\\
F_0 \cdot (\omega - 1) \cdot d_0{}^2 \cdot a^6 = 0 \enspace.
\end{cases}
\]
Now if $\omega + 2 = 0$ in $\FF_{p^2}$, then we have $\omega = -2$ and $1 = \omega^3 = -8$, therefore $9 = 0$, contradicting the fact $p \neq 3$.
Hence we have $\omega + 2 \neq 0$.
Similarly, if $2 \omega + 1 = 0$ in $\FF_{p^2}$, then we have $\omega = -1/2$ and $1 = \omega^3 = -1/8$, therefore $9 = 0$, contradicting the fact $p \neq 3$.
Hence we have $2 \omega + 1 \neq 0$.
By this and the facts $\omega \neq 1$, $p \not\in \{2,3\}$, $d_0 \neq 0$, and $a \neq 0$, it follows that
\[
F_2 \cdot 5 \cdot 7 \cdot 13 = 0 \,,\,
F_1 = 0 \,,\,
F_0 = 0 \enspace.
\]
We divide the proof into the following four cases.

\paragraph{Case 1: $p \not\in \{5,7,13\}$.}

In this case, we have $F_2 = F_1 = F_0 = 0$.
Now $F_2$ can be factored as
\[
F_2 = (z + 8) \cdot (z^2 - 20 z - 8) \cdot (z^2 + 2626 z - 656) \enspace,
\]
therefore we have $z + 8 = 0$, $z^2 - 20 z - 8 = 0$, or $z^2 + 2626 z - 656 = 0$.
We divide the proof into the following three cases.

\paragraph{Case 1-1: $z + 8 = 0$.}

Now the remainder of the polynomial $F_0$ modulo $z + 8$ is $1632586752000 = 2^{13} \cdot 3^{13} \cdot 5^3$, which must be zero as $F_0 = 0$ and $z + 8 = 0$.
This contradicts the assumption $p \not\in \{2,3,5\}$.

\paragraph{Case 1-2: $z^2 - 20 z - 8 = 0$.}

Now the remainders of the polynomials $F_0$ and $F_1$ modulo $z^2 - 20 z - 8$ are
\[
-2714622253056 z - 1066458742272
= -2^9 \cdot 3^8 \cdot 7^2 \cdot 19 \cdot 31 \cdot (28 z + 11)
\]
and
\[
3708367899264 z + 1454261921280
= 2^7 \cdot 3^9 \cdot 7^2 \cdot 19 \cdot 31 \cdot (51 z + 20) \enspace,
\]
respectively.
These two polynomials must be zero as $F_0 = F_1 = 0$ and $z^2 - 20 z - 8 = 0$.
We divide the proof into the following two cases.

\paragraph{Case 1-2-1: $p \not\in \{19,31\}$.}

In this case, by the assumption $p \not\in \{2,3,7\}$, we have $28z + 11 = 51z + 20 = 0$.
Hence
\[
0 = 51 \cdot (28 z + 11) - 28 \cdot (51 z + 20)
= 1 \enspace,
\]
a contradiction.

\paragraph{Case 1-2-2: $p \in \{19,31\}$.}

In this case, for any choice of $p$, we have $z^{(p^2 - 1) / 3} \bmod (z^2 - 20z - 8) = 1$ as polynomials in $z$ over $\FF_p$, therefore the element $z \in \FF_{p^2}$ satisfying $z^2 - 20z - 8 = 0$ also satisfies $z^{(p^2 - 1) / 3} = 1$.
Hence $z = d_0{}^3 / a$ is a cube in $\FF_{p^2}$.
This implies that $a$ is also a cube in $\FF_{p^2}$, contradicting the assumption that $c \not\in \FF_{p^2}$.

\paragraph{}

Hence we have a contradiction in all subcases for Case 1-2.

\paragraph{Case 1-3: $z^2 + 2626 z - 656 = 0$.}

Now the remainders of the polynomials $F_0$ and $F_1$ modulo $z^2 + 2626 z - 656$ are
\[
31682886407915875200 z - 7913936287671782400
= 2^7 \cdot 3^8 \cdot 5^2 \cdot 37 \cdot 53 \cdot (769532791 z - 192218392)
\]
and
\[
124498590050861874000 z - 31097984477624784000
= 2^4 \cdot 3^8 \cdot 5^3 \cdot 37 \cdot 53 \cdot (4838233297 z - 1208522152) \enspace,
\]
respectively.
These two polynomials must be zero as $F_0 = F_1 = 0$ and $z^2 + 2626 z - 656 = 0$.
We divide the proof into the following two cases.

\paragraph{Case 1-3-1: $p \not\in \{11,17,23,29,37,53\}$.}

In this case, by the assumption $p \not\in \{2,3,5\}$, we have
\[
769532791 z - 192218392
= 4838233297 z - 1208522152
= 0 \enspace.
\]
Hence
\[
\begin{split}
0 &= 4838233297 \cdot (769532791 z - 192218392) - 769532791 \cdot (4838233297 z - 1208522152) \\
&= 143687808
= 2^7 \cdot 3^2 \cdot 11 \cdot 17 \cdot 23 \cdot 29 \enspace,
\end{split}
\]
contradicting the assumption $p \not\in \{2,3,11,17,23,29\}$.

\paragraph{Case 1-3-2: $p \in \{11,17,23,29,37,53\}$.}

In this case, similarly to Case 1-2-2, for any choice of $p$, we have $z^{(p^2 - 1) / 3} \bmod (z^2 + 2626 z - 656) = 1$ as polynomials in $z$ over $\FF_p$, therefore the element $z \in \FF_{p^2}$ satisfying $z^2 + 2626 z - 656 = 0$ also satisfies $z^{(p^2 - 1) / 3} = 1$.
This implies a contradiction in a way similar to Case 1-2-2.

\paragraph{}

Hence we have a contradiction in all subcases for Case 1-3, therefore we have a contradiction in all subcases for Case 1.

\paragraph{Case 2: $p = 5$.}

In this case, we have
\[
\gcd(F_0,F_1)
= z^5 + 2 z^4 + 4 z^3 + 3 z^2 + z
= z (z + 3)^4
\]
as polynomials in $z$ over $\FF_p$, therefore $z (z + 3)^4 = 0$ as $F_1 = F_0 = 0$.
This implies that $z = -3$ (as $z = d_0{}^3 / a \neq 0$), which is a cube in $\FF_p$ (as now $p \equiv 2 \pmod{3}$).
This implies a contradiction in a way similar to Case 1-2-2.

\paragraph{Case 3: $p = 7$.}

In this case, we have
\[
\gcd(F_0,F_1)
= z^4 + 2 z^3 + 6 z^2 + 5 z + 1
= (z^2 + z + 6)^2
\]
as polynomials in $z$ over $\FF_p$, therefore $(z^2 + z + 6)^2 = 0$ as $F_1 = F_0 = 0$, hence $z^2 + z + 6 = 0$.
Now we have $z^{(p^2 - 1) / 3} \bmod (z^2 + z + 6) = 1$ as polynomials in $z$ over $\FF_p$.
This implies a contradiction in a way similar to Case 1-2-2.

\paragraph{Case 4: $p = 13$.}

In this case, we have
\[
\gcd(F_0,F_1)
= z^4 + 6 z^3 + 12 z^2 + 3 z + 5
= (z^2 + z + 12) (z^2 + 5 z + 8)
\]
as polynomials in $z$ over $\FF_p$, therefore $z^2 + z + 12 = 0$ or $z^2 + 5 z + 8 = 0$ as $F_1 = F_0 = 0$.
Now we have $z^{(p^2 - 1) / 3} \bmod (z^2 + z + 12) = 1$ and $z^{(p^2 - 1) / 3} \bmod (z^2 + 5 z + 8) = 1$ as polynomials in $z$ over $\FF_p$.
This implies a contradiction in a way similar to Case 1-2-2.

\paragraph{}

Hence we have a contradiction in any case.
Therefore we have $c \in \FF_{p^2}$, as desired.
This concludes the proof of Theorem \ref{thm:whether_codomain_is_over_Fp2}.
\end{proof}

\end{document}
